\section{总结与延伸}

这堂嘉宾课最终把推理系统和模型研究统一起来。服务端不是训练完成后的附属工程，而是新的观测窗口：真实 workload 揭示 cache 与 scheduler 问题，GPU 时间线揭示 kernel 空洞，memory footprint 又提示参数复用和循环架构的价值。

\begin{figure}[H]
\centering
\includegraphics[width=0.9\textwidth]{01-00-00-takeaway.jpg}
\caption{讲者的收束：理解 inference 与 GPU kernels，才能做 full-stack ML innovation。}
\figsource{课堂视频 00:59:33--01:00:30。}
\end{figure}

\subsection{一张端到端设计检查表}

为了把整堂课变成可执行流程，下面的检查表按层次记录“先问什么”和“必须观测什么”。它的核心用法不是逐项打勾，而是沿着因果链追踪：workload 变化是否先改变 SLA，随后改变 queue/cache，最后才在 kernel 或 architecture 层表现出来。

{\small
\setlength{\tabcolsep}{4pt}
\begin{longtable}{L{0.18\textwidth}L{0.36\textwidth}L{0.36\textwidth}}
\toprule
层次 & 先问什么 & 需要记录什么 \\
\midrule
Workload & 输入/输出长度、会话轮数、到达间隔怎样分布？ & histogram、cache hit、cold/warm 比例、工具调用模式 \\
SLA & 用户在乎首 token、流畅度还是总体完成时间？ & p50/p95/p99 TTFT、TBT、完成率、QPS/GPU \\
Scheduler & 谁能进入 batch，谁被谁阻塞？ & queue delay、preemption、KV blocks、priority、admission \\
Cache & 哪些状态值得留在 HBM，何时 prefetch？ & block size、reuse distance、eviction reason、transfer time \\
Parallelism & 模型、expert、sequence 怎样切分？ & topology、collective time、failure domain、imbalance \\
Kernel & 时间线空洞来自 launch、load 还是 tail？ & SM utilization、bandwidth、kernel boundaries、shape \\
Architecture & 参数、recurrence、KV 大小如何匹配硬件？ & weight footprint、loop count、quantization、target hardware \\
Validation & 优化是否只对单一 benchmark 有效？ & mixed traffic replay、tail SLA、quality regression、故障演练 \\
\bottomrule
\caption{从 workload 到 architecture 的全栈审计表。}
\end{longtable}
}

\subsection{五条最终结论}

\begin{enumerate}
  \item \textbf{先定义 workload，再定义最优。} 没有输入/输出长度、会话和 cache 分布，“最快引擎”没有可比较含义。
  \item \textbf{KV cache 是在线推理的核心状态。} Continuous batching、prefix sharing、offload、CPD 和 agentic serving 都围绕它展开。
  \item \textbf{Decode 优化的核心常是搬运与调度。} Megakernel 通过消除 boundary、重叠 load/compute 接近硬件带宽上限。
  \item \textbf{系统约束可以产生架构创新。} Parcae 用 recurrence 增加计算而不同比例增加参数，并用动力系统约束稳定性。
  \item \textbf{Full-stack co-design 需要承认工程成本。} 专用 kernel、硬件格式与模型结构能带来巨大收益，但可移植性、验证和维护必须进入目标函数。
\end{enumerate}

\enlargethispage{2\baselineskip}
\subsection{拓展阅读}

\begin{itemize}
  \item Parcae 论文与代码：\href{https://arxiv.org/abs/2604.12946}{arXiv:2604.12946}，\href{https://github.com/sandyresearch/parcae}{SandyResearch/parcae}；
  \item Cache-aware prefill--decode disaggregation：\href{https://www.together.ai/blog/cache-aware-disaggregated-inference}{Together AI 官方说明}；
  \item Whole-model megakernel：\href{https://hazyresearch.stanford.edu/blog/2025-05-27-no-bubbles}{HazyResearch: Look Ma, No Bubbles!}；
  \item Kernel 框架：\href{https://github.com/HazyResearch/ThunderKittens}{HazyResearch/ThunderKittens} 与 \href{https://github.com/HazyResearch/Megakernels}{HazyResearch/Megakernels}。
\end{itemize}

\end{document}
